% Part: methods
% Chapter: proofs
% Section: cant-do-it

\documentclass[../../../include/open-logic-section]{subfiles}

\begin{document}

\olfileid{mth}{prf}{cnt}

\olsection{“我做不到！”}

我们都有感到想要放弃的时刻。但你\emph{确实}能做到。你的授课教师和助教，以及你的同学们，都可以提供帮助。去向他们求助吧！这里有一些小贴士，可以帮助你避免崩溃，以及在你想要放弃时该怎么做。

为了确保你能顺利解决问题，请做到以下几点：

\begin{enumerate}
\item \emph{尽可能早地开始。}整个学期我们都会很忙，而且我们很多人都在与拖延症抗争，你能做的最有效的事情之一就是尽早开始做作业。这样，如果你卡壳了，你还有时间去寻找解决办法（而不是崩溃大哭）。
\item \emph{和你的同学交流。}你并不孤单。班上的其他人可能也会遇到困难——但他们遇到的困难可能与你不同。与同伴讨论可以让你从不同的角度看待问题，这可能会带来突破。当然，不要直接抄他们的答案：向他们请求提示，或者解释你卡在哪儿了，请他们告诉你下一步该怎么做。而当你自己弄明白之后，也要回馈他们。帮助别人、向别人解释，也会让你自己理解得更透彻。
\item \emph{主动求助。}你有很多可用的资源——你的授课教师和助教就在那里为你提供帮助，并且\emph{确实希望}你成功。他们应该能帮助你解决问题，并找出你在哪个环节遇到了困难。
\item \emph{休息一下。}如果你卡壳了，那\emph{可能}是因为你盯着这个问题太久了。短暂休息一下，喝杯茶，或者先做一会儿别的题目，然后再以清醒的头脑回到这个问题上。睡一觉再说。
\end{enumerate}

注意到了吗？这些策略都要求你尽早着手做证明。如果你在截止日期前一天的凌晨两点才开始写证明，这些策略可能就没那么有用了。

这听起来可能有点悲观，但攻克一道证明题是一项最终会有回报的挑战。有些人以此为职业——所以其中肯定有让人乐在其中的地方。和几乎所有事情一样，解决问题和做证明是需要练习的。你可能会看到一些同学觉得这很容易：他们很可能只是已经经过了大量练习。尽量不要轻易放弃。

如果你在某个具体问题上确实耗尽了时间（或耐心）：那也没关系。这并不意味着你笨，也不意味着你永远都搞不定。去向你的老师或其他同学请教正确的做法，并找出你在哪里出错了或卡住了，这样下次遇到类似问题时你就可以避免重蹈覆辙。然后，试着不看答案自己独立做一遍。下次，记得早点开始（早点求助）。

\end{document}